%------------------------------------------------------------------------------%

\usepackage{integracao_e_series}
\usepackage{systeme}

\title{Integração por Frações Parciais}
\subtitle{Caso 3 -- Raizes complexas}

%------------------------------------------------------------------------------%
\begin{document}

\addtitle

%------------------------------------------------------------------------------%
\section{Integração por Frações Parciais}

%------------------------------------------------------------------------------%
\begin{frame}{Caso 3}
  $Q(x)$ contém fatores quadráticos irredutíveis que não se repetem
  \[
    Q(x) = \cdots  \left( ax^2 + bx + c \right)
    \qquad
    \text{ com } b^2-4ac<0
  \]

  \pause
  Expansão
  \[
    \frac{P(x)}{Q(x)}
    = \cdots
    + \frac{Ax+B}{ax^2 + bx + c}
  \]
\end{frame}
%------------------------------------------------------------------------------%

\begin{frame}{Exemplo}
  \[
    \frac{x}{(x-2)(x^2+1)(x^2+4)}
    = \frac{A}{x-2}
    + \frac{Bx+C}{x^2+1}
    + \frac{Dx+E}{x^2+4}
  \]
\end{frame}

% 3
%------------------------------------------------------------------------------%
\addtocounter{exemplo}{1}
\section{Exemplo \theexemplo}

%------------------------------------------------------------------------------%
\begin{frame}{Exemplo \theexemplo}
  Encontre \(\int \frac{2x^2-x+4}{x^3+4x}\,dx\)

  \pause
  O grau do numerador $P$ é menor que o grau do denominador $Q$

  \pause
  Fatorando $Q$
  \[
    Q(x)
    = x^3 + 4x
    = x\left( x^2 + 4 \right)
  \]
  \pause
  Caso 3: Fator quadrático irredutível sem repetição
\end{frame}

%------------------------------------------------------------------------------%
\begin{frame}{Exemplo \theexemplo}
  \onslide<+->{\[
    \frac{ 2x^2 - x + 4}{ x^3 + 4x }
    = \frac{2x^2 - x + 4}{x\left( x^2 + 4 \right)}
    = \frac{A}{x} + \frac{Bx+C}{x^2+4}
  \]}
  \onslide<+->{Multiplicando por \(x^3 + 4x=x\left(x^2+4\right)\)}
  \begin{align*}
    \onslide<+->{2x^2 - x + 4
    & = A\left( x^2 + 4 \right) + \left( Bx + C \right) x \\}
    \onslide<+->{& = \left( A + B \right) x^2 + Cx + 4A}
  \end{align*}
  \onslide<+->{\systeme[ABC]{
    A+B=2,
    C=-1,
    4A=4
    }}
  \onslide<+->{\[
    A =  1 \qquad
    B =  1 \qquad
    C = -1
  \]}
\end{frame}

%------------------------------------------------------------------------------%
\begin{frame}{Exemplo \theexemplo}
  \begin{align*}
    \onslide<+->{\int \frac{2x^2-x+4}{x^3+4x}dx
    & = \int \frac{1}{x} + \frac{x-1}{x^2+4} \,dx \\[3mm]}
    \onslide<+->{& = \int \frac{1}{x}dx + \int \frac{x}{x^2+4} dx - \int \frac{1}{x^2+4} dx \\[3mm]}
    \onslide<+->{& = \ln\abs{x}
    + \frac{1}{2} \ln\abs{x^2+4}
    - \frac{1}{2} \arctan\left(\frac{x}{2}\right)
    + C}
  \end{align*}
  \onslide<+->{Onde usamos
  \[
    u=x^2+4 \qquad \text{para calcular} \qquad\int \frac{x}{x^2+4} dx
  \]}
  \onslide<+->{\[
    \int \frac{1}{x^2+a^2}dx
    = \frac{1}{a} \arctan\left(\frac{x}{a}\right)+K
  \]}
\end{frame}

%------------------------------------------------------------------------------%
\section{Lista Mínima}

%------------------------------------------------------------------------------%
\begin{frame}{Lista Mínima}
  Estudar a Seção 4.5 da Apostila

  Exercícios:

  \vfill
  Atenção: A prova é baseada no livro, não nas apresentações
\end{frame}

%------------------------------------------------------------------------------%
\end{document}
%------------------------------------------------------------------------------%